\section{Relation to Established Optimization Methods}\label{sec:r37-position}
\paragraph{Contracts and retained information.}
Continuation promises and participation are established state variables in dynamic contracting \citep{SpearSrivastava1987,ThomasWorrall1988}. Finite policy graphs also have a contractual history \citep{Zhang2012}. Our model has public exogenous branches and no private reports; it is not an incentive-compatibility result for hidden types. Its resource constraint is the information crossing a specified execution boundary. Classical behavioral minimization determines exact equivalence once a finite output transducer has been constructed \citep{GivanDeanGreig2003,Peyriere2023}. Here the design problem deliberately operates below the exact alphabet size.

The companion maps the unfolded accepted-service problem invertibly to laminar resource allocation. Tree-constrained allocation, polymatroid optimization, and normalized memoization are inherited tools \citep{Mjelde1983,Tang1990,FedergruenGroenevelt1986,Groenevelt1991}. The compact-input event and storage bounds concern a repeated-subtree subclass, not a new solution to arbitrary laminar allocation. Those results support the resource accounting but are not needed to solve the renewal frontier developed below.

\paragraph{Four algorithmic layers.}
Quadrangle-inequality dynamic programming is classical \citep{Yao1980}. Searching a full totally monotone matrix is the SMAWK result of \citet{Aggarwal1987}. Partial matrices require a separate comparison: \citet{KlaweKleitman1990} developed almost-linear generalized searching, and \citet{Chan2021} obtained expected linear time for the harder Monge staircase setting. Submatrix-query data structures for full and partial Monge matrices form a further related literature \citep{Kaplan2017}. These problems differ from a single completed row-minimum pass; citing only full-matrix SMAWK would obscure the relevant prior work.

Most importantly, the introduction of \citet[pp.~2--3 of the author version]{Chan2021} already separates staircase orientations that admit direct completion and SMAWK from the harder orientations. Our terminal array has feasible suffixes that shrink with increasing row index. Under the row-minimum convention used below, this is the concave-Monge falling-staircase easy case. The program array has feasible prefixes that expand. Reversing both its row and column orders converts it to the same easy case and preserves the Monge inequality. We give direct prefix and suffix proofs so that the actual left-tie convention, rather than an implicit reversal of tie order, is audited. Neither proof claims a new general completion method.

The distinction also explains the complexity comparison. No generic inverse-Ackermann barrier is overcome here. The branch order fixes a favorable nesting direction; completed row preferences therefore cannot reverse in the direction forbidden by total monotonicity. Other partial domains, the opposite orientation under a fixed inequality convention, and online dependencies within a searched layer are outside our theorem. All predecessor values in a layer are known before its search begins. The quantitative contractual work is to establish the finite-cost inequalities after continuous minimization and a constant-work exact oracle. Once those facts hold, the linear search is a classical consequence.

\citet{Emmerich2026} provides a recent, different application: integral R2 subset selection is reduced to a Monge dynamic program and accelerated through monotone and staircase search. That work reinforces that this generic algorithmic progression is not distinctive to service contracts. It does not supply the present cap-anchoring reduction, institutional comparison, or continuously minimized contractual terminal costs. Table~\ref{tab:r37-priority} separates these contributions theorem by theorem. We attribute the search machinery positively rather than assign novelty to a change in implementation order.
